% Generated by publication/build_sections.py; do not edit by hand.
\paragraph{Informal verdict.}
\textbf{A --- Complete solution}, with high confidence.

\paragraph{Lean coverage (partial).}
Lean proves the exact minimum and complete block-median iff classification for a supplied integrable nondecreasing rearrangement q, permutation invariance/sorting, continuous positive-block median uniqueness, and Phi\_eq\_of\_map\_eq conditional on a supplied equal-pushforward q. The sole canonical gap is to construct such a q from every printed continuous integrable f and expose the unconditional original-f minimum and all-minimizers iff theorem. Strictness, asymptotics, Bellman/stationarity refinements, affine formulas, compactness refinements, smooth nonuniqueness, and a separately named constant-function corollary are optional and are not follow-up completion requirements.

\subsection{Audited informal answer}
Q766 --- Complete solution
\begin{enumerate}
\item Faithful restatement
\end{enumerate}
Let \(f:(0,1)\to\mathbb R\) be continuous and integrable, and let \(n\ge 1\). For \(X=(x_1,\dots,x_n)\in\mathbb R^n\), define

\[
\delta(X,y)=\min_{1\le k\le n}|y-x_k|,
\qquad
\phi(X)=\int_0^1\delta(X,f(t))\,dt.
\]

The problem asks for the minimum of \(\phi\) and for the complete set of \(n\)-tuples at which it is attained. The statement explicitly permits adding assumptions on \(f\) when useful and notes that the case \(n=1\) was treated in an earlier examination problem. There are no subparts, footnotes, or apparent typographical errors. Luc Pommeret
I interpret ``integrable'' in the standard French sense

\[
\int_0^1 |f(t)|\,dt<\infty.
\]

This is the natural necessary condition: if only a conditionally convergent improper integral were meant and \(\int_0^1|f|=\infty\), then \(\phi(X)=+\infty\) for every finite tuple \(X\).

\begin{enumerate}
\item Distribution and increasing rearrangement
\end{enumerate}
Let \(\lambda\) denote Lebesgue measure on \((0,1)\), and let

\[
\mu=f_\#\lambda
\]

be the image probability measure of \(f\); explicitly,

\[
\mu(B)=\lambda\bigl(\{t\in(0,1):f(t)\in B\}\bigr).
\]

Then

\[
\phi(x_1,\dots,x_n)
=
\int_{\mathbb R}\min_{1\le k\le n}|y-x_k|\,d\mu(y).
\]

Thus the problem depends only on the distribution of \(f\), not on the order in which its values occur.
Let

\[
F(y)=\mu((-\infty,y])
\]

and define the generalized quantile, or nondecreasing rearrangement, by

\[
q(u)=\inf\{y\in\mathbb R:F(y)\ge u\},
\qquad 0<u<1.
\]

Then \(q\) is nondecreasing, \(q\in L^1(0,1)\), and \(q_\#\lambda=\mu\). Hence

\[
\phi(X)=\int_0^1\min_{1\le k\le n}|q(u)-x_k|\,du.
\]

Set

\[
H(s)=\int_0^s q(u)\,du,\qquad 0\le s\le1.
\]

The function \(H\) is continuous on \([0,1]\) and convex.
For \(0\le a\le b\le1\), define

\[
K(a,b)
=
H(a)+H(b)-2H\!\left(\frac{a+b}{2}\right).
\]


\begin{enumerate}
\item Main theorem: exact minimum and complete classification
\end{enumerate}
Let

\[
\Delta_n
=
\left\{(s_0,\dots,s_n):
0=s_0\le s_1\le\cdots\le s_n=1
\right\}
\]

and, for \(s\in\Delta_n\), put

\[
J_n(s)
=
\sum_{i=1}^n
K(s_{i-1},s_i).
\]

Theorem 1
For every integrable measurable \(f:(0,1)\to\mathbb R\),

\[
\boxed{
\min_{X\in\mathbb R^n}\phi(X)
=
\min_{s\in\Delta_n}
\sum_{i=1}^n
\left[
H(s_{i-1})+H(s_i)
-2H\!\left(\frac{s_{i-1}+s_i}{2}\right)
\right].
}
\tag{3.1}
\]

Both minima are attained.
More precisely, let

\[
\mathcal S_n=\operatorname*{argmin}_{s\in\Delta_n}J_n(s).
\]

For \(s\in\mathcal S_n\) and every positive-length block
\(s_{i-1}<s_i\), set

\[
\theta_i=\frac{s_{i-1}+s_i}{2}
\]

and define

\[
M_i(s)=
[q(\theta_i-),q(\theta_i+)],
\]

where \(q(\theta_i-)\) and \(q(\theta_i+)\) are the one-sided limits of the monotone function \(q\).
Choose any

\[
c_i\in M_i(s)
\]

on every positive block; on a zero-length block one may choose \(c_i\) arbitrarily. Then

\[
(c_1,\dots,c_n)
\]

is a minimizer of \(\phi\).
Conversely, every minimizing tuple, after reordering its coordinates increasingly, arises in this way from an element of \(\mathcal S_n\).
Thus (3.1), together with the median intervals \(M_i(s)\), is a complete description of both the minimum and its entire attainment set.

\begin{enumerate}
\item Proof of Theorem 1
\end{enumerate}
4.1 The one-centre problem
We first recall and prove the elementary median principle.
Let \(\nu\) be a finite positive measure of total mass \(p\), with finite first moment, and let

\[
G(c)=\int_{\mathbb R}|y-c|\,d\nu(y).
\]

The function \(G\) is convex. Its one-sided derivatives are

\[
G'_-(c)
=
\nu((-\infty,c))-\nu([c,\infty)),
\]

and

\[
G'_+(c)
=
\nu((-\infty,c])-\nu((c,\infty)).
\]

Indeed, these formulas follow by applying dominated convergence to the one-sided difference quotients of \(|y-c|\).
A point \(c\) minimizes \(G\) if and only if

\[
G'_-(c)\le0\le G'_+(c),
\]

equivalently,

\[
\nu((-\infty,c))\le \frac p2,
\qquad
\nu((c,\infty))\le \frac p2.
\tag{4.1}
\]

These are exactly the median inequalities.
Now restrict \(q\) to a quantile interval \((a,b)\), where \(a<b\), and let

\[
m=\frac{a+b}{2}.
\]

Since \(q\) is nondecreasing, the complete median interval of \(q\) on \((a,b)\) is

\[
[q(m-),q(m+)].
\]

For any \(c\) in this interval,

\[
q(u)\le c\quad\text{for a.e. }u\in(a,m),
\]

and

\[
q(u)\ge c\quad\text{for a.e. }u\in(m,b).
\]

Because both halves have the same length,

\[
\begin{aligned}
\min_{c\in\mathbb R}\int_a^b|q(u)-c|\,du
&=
\int_a^m(c-q(u))\,du
+
\int_m^b(q(u)-c)\,du\\
&=
\int_m^b q(u)\,du-\int_a^m q(u)\,du\\
&=
H(a)+H(b)-2H(m)\\
&=K(a,b).
\end{aligned}
\tag{4.2}
\]

For \(a=b\), both sides are \(0\).

4.2 From centres to an ordered quantile partition
Fix \(X=(x_1,\dots,x_n)\). Since \(\phi\) is invariant under permutations, assume

\[
x_1\le x_2\le\cdots\le x_n.
\]

For each \(y\in\mathbb R\), choose the smallest index realizing the nearest-centre distance:

\[
r(y)=
\min\left\{
i: |y-x_i|=\min_j|y-x_j|
\right\}.
\]

The map \(r\) is nondecreasing.
Indeed, suppose \(y<z\), but \(r(y)=j>i=r(z)\). Then \(x_i\le x_j\). The inequality saying that \(x_j\) is at least as close to \(y\) as \(x_i\) gives

\[
y\ge\frac{x_i+x_j}{2},
\]

while the inequality saying that \(x_i\) is at least as close to \(z\) as \(x_j\) gives

\[
z\le\frac{x_i+x_j}{2},
\]

contradicting \(y<z\).
Since both \(q\) and \(r\) are nondecreasing, \(r\circ q\) is a nondecreasing integer-valued function on \((0,1)\). Consequently its level sets form, up to endpoints, consecutive intervals

\[
(s_{i-1},s_i),\qquad
0=s_0\le s_1\le\cdots\le s_n=1.
\]

Using (4.2),

\[
\begin{aligned}
J_n(s)
&=
\sum_{i=1}^n
\min_{c\in\mathbb R}
\int_{s_{i-1}}^{s_i}|q(u)-c|\,du\\
&\le
\sum_{i=1}^n
\int_{s_{i-1}}^{s_i}|q(u)-x_i|\,du\\
&=
\int_0^1\min_j|q(u)-x_j|\,du\\
&=\phi(X).
\end{aligned}
\]

Therefore

\[
\min_{s\in\Delta_n}J_n(s)
\le
\inf_{X\in\mathbb R^n}\phi(X).
\tag{4.3}
\]


4.3 From a quantile partition to centres
Conversely, take any \(s\in\Delta_n\). For every positive block choose a median \(c_i\in M_i(s)\); on empty blocks choose anything. Then, pointwise,

\[
\min_j|q(u)-c_j|
\le |q(u)-c_i|,
\qquad s_{i-1}<u<s_i.
\]

Hence

\[
\begin{aligned}
\phi(c_1,\dots,c_n)
&=
\int_0^1\min_j|q(u)-c_j|\,du\\
&\le
\sum_{i=1}^n
\int_{s_{i-1}}^{s_i}|q(u)-c_i|\,du\\
&=
J_n(s).
\end{aligned}
\]

Taking the infimum over \(X\) and then the minimum over \(s\),

\[
\inf_X\phi(X)
\le
\min_{s\in\Delta_n}J_n(s).
\tag{4.4}
\]

Combining (4.3) and (4.4) proves equality.
Since \(H\) is continuous, \(J_n\) is continuous on the compact simplex \(\Delta_n\). Thus the right-hand minimum is attained. Choosing medians on an optimal partition gives a tuple attaining the same value, so the left-hand minimum is also attained.

4.4 Equality cases
Suppose first that \(s\in\mathcal S_n\) and the \(c_i\) are block medians. We obtained

\[
\phi(c_1,\dots,c_n)\le J_n(s)=\min_X\phi(X).
\]

The reverse inequality follows from the definition of the minimum, so equality holds: every such tuple is minimizing.
Conversely, let \(X\) be minimizing and form its nearest-centre quantile partition \(s\) as in §4.2. Then

\[
\min_sJ_n(s)
\le
J_n(s)
\le
\phi(X)
=
\min_X\phi(X)
=
\min_sJ_n(s).
\]

All inequalities are equalities. In particular \(s\in\mathcal S_n\), and on each nonempty cell

\[
\int_{s_{i-1}}^{s_i}|q(u)-x_i|\,du
=
\min_c\int_{s_{i-1}}^{s_i}|q(u)-c|\,du.
\]

Therefore \(x_i\) is a median on that block, i.e.

\[
x_i\in M_i(s).
\]

This proves the full classification.

\begin{enumerate}
\item Specialization to the continuous function in Q766
\end{enumerate}
Continuity has a useful additional consequence.
Lemma 2: the distribution of a nonconstant continuous \(f\) has no gaps
Let

\[
I=f((0,1)).
\]

Since \((0,1)\) is connected and \(f\) is continuous, \(I\) is an interval.
Moreover,

\[
\operatorname{supp}\mu=\overline I.
\]

Indeed, if \(y=f(t_0)\) and \(U\) is a neighbourhood of \(y\), continuity gives an interval \(V\ni t_0\) such that \(f(V)\subset U\). Thus

\[
\mu(U)\ge\lambda(V)>0.
\]

The same argument applies to points in \(\overline I\) by first choosing a value of \(f\) inside the prescribed neighbourhood.
Consequently, if \(f\) is nonconstant, \(\operatorname{supp}\mu\) is a nondegenerate interval, possibly unbounded. There are no open gaps of zero \(\mu\)-measure inside the support.
A jump of the quantile \(q\) would correspond to such a gap. Hence:

\[
\boxed{f\text{ nonconstant and continuous }\Longrightarrow q
\text{ is continuous on }(0,1).}
\tag{5.1}
\]

The distribution may have atoms---\(f\) can be constant on an interval---but atoms produce flat portions of \(q\), not jumps.
Therefore every positive quantile block has a unique median:

\[
M_i(s)=\{q(\theta_i)\}.
\]


Theorem 3: complete answer for the original continuous \(f\)
Constant case
If \(f(t)=c\) for all \(t\), then

\[
\phi(X)=\min_k|c-x_k|.
\]

Thus

\[
\boxed{\min\phi=0,}
\]

and

\[
\boxed{
\operatorname{Argmin}\phi
=
\{X\in\mathbb R^n:\text{some coordinate of }X\text{ equals }c\}.
}
\]

The other coordinates are arbitrary.
Nonconstant case
Assume \(f\) is nonconstant. Let \(q\), \(H\), \(J_n\), and \(\mathcal S_n\) be as above. Then

\[
\boxed{
\min_{X\in\mathbb R^n}\phi(X)
=
\min_{0=s_0\le\cdots\le s_n=1}
\sum_{i=1}^n
\left[
H(s_{i-1})+H(s_i)
-2H\!\left(\frac{s_{i-1}+s_i}{2}\right)
\right].
}
\tag{5.2}
\]

Every optimal partition is in fact strict:

\[
0=s_0<s_1<\cdots<s_{n-1}<s_n=1.
\tag{5.3}
\]

The complete attainment set is

\[
\boxed{
\operatorname{Argmin}\phi
=
\bigcup_{s\in\mathcal S_n}
\left\{
\bigl(
q(\tfrac{s_{\sigma(1)-1}+s_{\sigma(1)}}2),
\dots,
q(\tfrac{s_{\sigma(n)-1}+s_{\sigma(n)}}2)
\bigr):
\sigma\in\mathfrak S_n
\right\}.
}
\tag{5.4}
\]

Equivalently, for an increasing ordering of the coordinates,

\[
\boxed{
x_i=q\!\left(\frac{s_{i-1}+s_i}{2}\right),
\qquad 1\le i\le n,
}
\tag{5.5}
\]

for some \(s\in\mathcal S_n\).
All optimal coordinates are distinct, and the minimizer set is compact.

Proof of strictness and distinctness
Write

\[
e_n=\min_{X\in\mathbb R^n}\phi(X).
\]

For nonconstant continuous \(f\),

\[
e_n<e_{n-1}\qquad(n\ge2).
\tag{5.6}
\]

To prove this, take an optimal \((n-1)\)-tuple \(A\). Its set of centres is finite, whereas \(\operatorname{supp}\mu\) is a nondegenerate interval. Choose

\[
y_0\in\operatorname{supp}\mu\setminus A
\]

and let

\[
d=\operatorname{dist}(y_0,A)>0.
\]

Choose \(0<r<d/3\). Since \(y_0\) lies in the support,

\[
\mu((y_0-r,y_0+r))>0.
\]

After adding \(y_0\) as a new centre, the distance decreases on this interval by at least

\[
(d-r)-r=d-2r>0
\]

and does not increase elsewhere. This proves (5.6).
If an optimal quantile partition for \(e_n\) had a zero-length block, deleting it would give an \((n-1)\)-block partition of the same cost, contradicting (5.6). This proves (5.3).
If two optimal centres coincided, the tuple would effectively use at most \(n-1\) distinct centres, again contradicting (5.6). Thus all centres are distinct.
Finally, \(\mathcal S_n\) is a compact subset of \(\Delta_n\). Since it does not meet the boundary where a block has zero length, all its block lengths are uniformly bounded below. Hence all the midpoint parameters stay in a compact subinterval of \((0,1)\), on which \(q\) is bounded and continuous. Formula (5.4) then shows that the minimizer set is compact.

\begin{enumerate}
\item Necessary equations for an optimum
\end{enumerate}
Let \(f\) be nonconstant and continuous, and let

\[
s\in\mathcal S_n,\qquad
\theta_i=\frac{s_{i-1}+s_i}{2},
\qquad
x_i=q(\theta_i).
\]

Since \(q\) is continuous on \((0,1)\),

\[
H'(u)=q(u),\qquad 0<u<1.
\]

Every internal breakpoint \(s_i\) is an unconstrained interior variable. Differentiating \(J_n\) with respect to \(s_i\) gives

\[
\begin{aligned}
0
&=
\frac{\partial J_n}{\partial s_i}\\
&=
\left[
q(s_i)-q\!\left(\frac{s_{i-1}+s_i}{2}\right)
\right]
+
\left[
q(s_i)-q\!\left(\frac{s_i+s_{i+1}}2\right)
\right].
\end{aligned}
\]

Thus every optimum satisfies

\[
\boxed{
2q(s_i)=x_i+x_{i+1},
\qquad 1\le i<n.
}
\tag{6.1}
\]

In value space, \(q(s_i)\) is therefore exactly the midpoint between the adjacent optimal centres. This is the Voronoi boundary condition.
Together, (5.5) and (6.1) say:


each centre is the median of the distribution in its Voronoi cell;


each boundary is the arithmetic midpoint of the adjacent centres.


These are the one-dimensional \(L^1\) centroidal Voronoi equations. They are necessary, but in general not sufficient for global optimality: the finite-dimensional function \(J_n\) need not be convex.

Atomless form
Suppose additionally that \(\mu\) is atomless. Then \(F\) is continuous and strictly increasing on its support, and \(q=F^{-1}\).
Let

\[
b_0=-\infty,\qquad
b_i=\frac{x_i+x_{i+1}}2\quad(1\le i<n),
\qquad
b_n=+\infty.
\]

Equation (6.1) gives

\[
s_i=F(b_i).
\]

Since \(x_i=q((s_{i-1}+s_i)/2)\),

\[
F(x_i)=\frac{s_{i-1}+s_i}{2}.
\]

Therefore every ordered optimal tuple satisfies

\[
\boxed{
2F(x_i)=F(b_{i-1})+F(b_i),
\qquad 1\le i\le n,
}
\tag{6.2}
\]

with \(F(-\infty)=0\) and \(F(+\infty)=1\).
Equivalently,

\[
\mu((b_{i-1},x_i))
=
\mu((x_i,b_i)).
\]

Thus \(x_i\) bisects the probability mass of its Voronoi cell.

\begin{enumerate}
\item Exact dynamic-programming algorithm
\end{enumerate}
Formula (5.2) admits an exact Bellman recursion.
For \(t\in[0,1]\), define

\[
V_0(0)=0,\qquad V_0(t)=+\infty\quad(t>0),
\]

and for \(k\ge1\),

\[
\boxed{
V_k(t)
=
\min_{0\le r\le t}
\bigl(
V_{k-1}(r)+K(r,t)
\bigr).
}
\tag{7.1}
\]

Then

\[
\boxed{e_n=V_n(1).}
\tag{7.2}
\]

All optimal partitions can be recovered by backtracking all minimizers \(r\) in (7.1). Once the breakpoints are known, the centres are given by the block medians, and for the nonconstant continuous case by

\[
x_i=q\!\left(\frac{s_{i-1}+s_i}{2}\right).
\]

A grid discretization of (7.1) with \(M\) quantile grid points has the standard \(O(nM^2)\) complexity, before any convexity or Monge-type acceleration.

\begin{enumerate}
\item Boundary and test cases
\end{enumerate}
8.1 The case \(n=1\)
There is no partition variable:

\[
s_0=0,\qquad s_1=1.
\]

Thus

\[
e_1
=
H(1)-2H\!\left(\frac12\right)
=
\int_{1/2}^1q(u)\,du-\int_0^{1/2}q(u)\,du.
\]

The complete median interval is

\[
[q(1/2-),q(1/2+)].
\]

For the nonconstant continuous \(f\), \(q\) is continuous, so the unique minimizer is

\[
\boxed{x_1=q(1/2),}
\]

the unique median of the distribution of \(f\).

8.2 Affine functions
Let

\[
f(t)=\alpha t+\beta,\qquad \alpha\ne0.
\]

The distribution is uniform on an interval of length \(|\alpha|\). Writing its lower endpoint as

\[
A=\min(\beta,\beta+\alpha),
\]

the quantile is

\[
q(u)=A+|\alpha|u.
\]

A direct calculation gives

\[
K(a,b)=\frac{|\alpha|}{4}(b-a)^2.
\]

Consequently,

\[
e_n
=
\frac{|\alpha|}{4}
\min_{\substack{h_i\ge0\\\sum h_i=1}}
\sum_{i=1}^n h_i^2
=
\frac{|\alpha|}{4n}.
\]

Equality in Cauchy--Schwarz forces

\[
h_i=\frac1n.
\]

Thus the unique increasing optimal tuple is

\[
\boxed{
x_i
=
A+\frac{2i-1}{2n}|\alpha|,
\qquad 1\le i\le n,
}
\]

and the full attainment set consists of its permutations.

8.3 Decay as \(n\to\infty\)
For every integrable \(f\),

\[
e_n\longrightarrow0.
\]

Indeed, given \(\varepsilon>0\), choose \(R>0\) such that

\[
\int_{|y|>R}|y|\,d\mu(y)<\varepsilon.
\]

Place \(n\) equally spaced centres in \([-R,R]\), including both endpoints. For \(|y|\le R\), the nearest-centre distance is at most \(R/(n-1)\), while outside \([-R,R]\) it is at most \(|y|\). Therefore

\[
e_n
\le
\frac{R}{n-1}
+
\int_{|y|>R}|y|\,d\mu(y)
<
\frac{R}{n-1}+\varepsilon.
\]

For nonconstant continuous \(f\), the sequence is therefore strictly decreasing and tends to zero:

\[
0<e_n<e_{n-1},
\qquad
e_n\downarrow0.
\]


\begin{enumerate}
\item Smooth high-resolution asymptotic
\end{enumerate}
The permission in the statement to add assumptions allows a more explicit asymptotic answer.
Theorem 4
Assume that the quantile \(q\) extends to a \(C^1\) function on \([0,1]\) and that

\[
q'(u)>0\qquad(0\le u\le1).
\]

Then

\[
\boxed{
e_n
=
\frac1{4n}
\left(
\int_0^1\sqrt{q'(u)}\,du
\right)^2
(1+o(1)).
}
\tag{9.1}
\]

Proof
Put

\[
g=q',\qquad
m_0=\min_{[0,1]}g>0,
\qquad
M_0=\max_{[0,1]}g,
\]

and let \(\omega\) be the modulus of continuity of \(g\).
For a block \([a,b]\), write

\[
h=b-a,\qquad c=\frac{a+b}{2}.
\]

Since \(q(c)\) is the block median,

\[
\begin{aligned}
K(a,b)
&=
\int_0^{h/2}
\bigl(q(c+v)-q(c-v)\bigr)\,dv\\
&=
\int_0^{h/2}
\int_{c-v}^{c+v}g(r)\,dr\,dv.
\end{aligned}
\]

Hence

\[
\left\lvert 
K(a,b)-\frac{g(c)h^2}{4}
\right\rvert 
\le
\frac{\omega(h)h^2}{4}.
\tag{9.2}
\]

Let

\[
A=\int_0^1\sqrt{g(u)}\,du.
\]

For the upper bound, choose the breakpoints so that

\[
\int_{s_{i-1}}^{s_i}\sqrt{g(u)}\,du=\frac An.
\]

The largest block length is \(O(1/n)\). By (9.2), uniformly in \(i\),

\[
K(s_{i-1},s_i)
=
\frac{A^2}{4n^2}(1+o(1)).
\]

Summing gives

\[
e_n\le\frac{A^2}{4n}(1+o(1)).
\tag{9.3}
\]

For the lower bound, an equal-length partition gives

\[
e_n\le\frac{M_0}{4n}.
\]

On the other hand, every block of length \(h\) costs at least

\[
\frac{m_0h^2}{4}.
\]

Thus, for an optimal partition, its largest block length tends to zero. Let \(h_i=s_i-s_{i-1}\) and \(c_i=(s_{i-1}+s_i)/2\). From (9.2),

\[
e_n
\ge
\frac{1-o(1)}4
\sum_{i=1}^n g(c_i)h_i^2.
\]

Cauchy--Schwarz gives

\[
\sum_{i=1}^n g(c_i)h_i^2
\ge
\frac1n
\left(
\sum_{i=1}^n\sqrt{g(c_i)}\,h_i
\right)^2.
\]

The sum in parentheses is a Riemann sum for \(A\), because the mesh tends to zero. Therefore

\[
e_n\ge\frac{A^2}{4n}(1-o(1)).
\tag{9.4}
\]

Combining (9.3) and (9.4) proves (9.1).

Density form
If \(\mu\) has a positive continuous density \(p\) on a compact interval \([A,B]\), then

\[
q'(u)=\frac1{p(q(u))}.
\]

Changing variables \(u=F(y)\) gives

\[
\int_0^1\sqrt{q'(u)}\,du
=
\int_A^B\sqrt{p(y)}\,dy.
\]

Hence

\[
\boxed{
e_n
=
\frac1{4n}
\left(
\int_A^B\sqrt{p(y)}\,dy
\right)^2
(1+o(1)).
}
\tag{9.5}
\]

If \(f\in C^1([0,1])\) is strictly monotone with \(f'\) nowhere zero, then \(q\) is \(f\) or its reversal, and

\[
\boxed{
e_n
=
\frac1{4n}
\left(
\int_0^1\sqrt{|f'(t)|}\,dt
\right)^2
(1+o(1)).
}
\tag{9.6}
\]

For an affine \(f\), (9.6) gives the exact value \(|\alpha|/(4n)\).

\begin{enumerate}
\item Uniqueness is not guaranteed
\end{enumerate}
Even smooth strict monotonicity of \(f\) does not force a unique optimal \(n\)-point set.
Consider first the nondecreasing step function

\[
q_0(u)=
\begin{cases}
-1,&0<u<1/5,\\
0,&1/5\le u\le4/5,\\
1,&4/5<u<1.
\end{cases}
\]

For \(n=2\), denote by \(J_{q_0}(s)\) the cost of the partition \((0,s,1)\). Directly,

\[
J_{q_0}(1/5)=\frac15,
\qquad
J_{q_0}(1/2)=\frac25.
\]

For two quantiles \(q,\widetilde q\), their primitives satisfy

\[
\sup_t|H_q(t)-H_{\widetilde q}(t)|
\le
\|q-\widetilde q\|_{L^1},
\]

and consequently

\[
\sup_{0\le s\le1}
|J_q(s)-J_{\widetilde q}(s)|
\le
8\|q-\widetilde q\|_{L^1}.
\tag{10.1}
\]

Choose a \(C^\infty\), strictly increasing function \(q\) satisfying

\[
q(1-u)=-q(u)
\]

and

\[
\|q-q_0\|_{L^1}<\frac1{80}.
\]

Such a \(q\) is obtained, for example, by symmetric mollification of \(q_0\), followed by adding a sufficiently small positive multiple of \(2u-1\).
By (10.1),

\[
J_q(1/5)<\frac3{10}
<
J_q(1/2).
\]

Thus \(1/2\) is not an optimal breakpoint. But symmetry gives

\[
J_q(s)=J_q(1-s).
\]

Therefore every minimizing breakpoint \(s_*\ne1/2\) has a distinct minimizing partner \(1-s_*\). Since \(q\) is strictly increasing, the corresponding unordered pairs of centres are distinct.
Taking \(f=q\) gives a smooth strictly increasing example with at least two different optimal two-point sets. Hence the occurrence of the possibly non-singleton set \(\mathcal S_n\) in (5.4) is genuinely necessary.

\begin{enumerate}
\item Role and necessity of the hypotheses
\end{enumerate}
Absolute integrability. This is necessary for a finite problem. If

\[
\int_0^1|f(t)|\,dt=\infty,
\]

then for \(M=\max_i|x_i|\),

\[
\delta(X,f(t))\ge (|f(t)|-M)_+,
\]

so \(\phi(X)=+\infty\) for every \(X\).
Continuity. It is not needed for the general quantile formula or existence theorem. It is used to show that the support is an interval with no gaps, hence that \(q\) is continuous and every positive block has a unique median. Without continuity, median intervals can be nontrivial. For example, for the distribution

\[
\frac12\delta_{-1}+\frac12\delta_1,
\]

every \(x\in[-1,1]\) is a one-centre minimizer.
Nonconstancy. It is necessary for strict decrease of \(e_n\), strict positivity of block lengths, and distinctness of optimal centres.
Atomlessness. It is needed only for the simple equal-mass equations (6.2). The exact quantile formula remains valid in the presence of atoms.
Compactness of the image. No such assumption is required. The support may be unbounded; finite first moment suffices.

\begin{enumerate}
\item Topological, countability, and choice assumptions
\end{enumerate}
The proof takes place entirely in the ordinary Euclidean spaces \(\mathbb R\), \(\mathbb R^n\), and the finite-dimensional simplex \(\Delta_n\).


The separation assumptions are only the usual Hausdorff structure of Euclidean space.


Existence uses the compactness of the closed simplex \(\Delta_n\) and continuity of \(J_n\). No compactness of \(\operatorname{supp}\mu\) is used.


All spaces involved are metrizable and second countable. Compactness may equivalently be handled by sequential compactness; nets or filters are unnecessary.


No countability assumption on level sets or fibres of \(f\) is used.


The quantile is canonically defined, so no choice principle is required there. Choosing one median on each of finitely many blocks uses only finite choice.


No interchange of a non-absolutely convergent integral or unjustified limiting operation occurs.



Final answer
For a nonconstant continuous integrable \(f\), let \(q\) be its nondecreasing rearrangement and \(H(s)=\int_0^s q\). Then

\[
\boxed{
\min_X\phi(X)
=
\min_{0=s_0<\cdots<s_n=1}
\sum_{i=1}^n
\left[
H(s_{i-1})+H(s_i)
-2H\!\left(\frac{s_{i-1}+s_i}{2}\right)
\right].
}
\]

The minimum is attained. If \(\mathcal S_n\) denotes the set of optimal partitions, the complete set of minimizing tuples is

\[
\boxed{
\bigcup_{s\in\mathcal S_n}
\mathfrak S_n\cdot
\left(
q\!\left(\frac{s_0+s_1}{2}\right),
\dots,
q\!\left(\frac{s_{n-1}+s_n}{2}\right)
\right).
}
\]

Every optimal partition satisfies

\[
\boxed{
2q(s_i)
=
q\!\left(\frac{s_{i-1}+s_i}{2}\right)
+
q\!\left(\frac{s_i+s_{i+1}}2\right),
\qquad 1\le i<n.
}
\]

These formulas give the exact minimum, a complete classification of every point of attainment, an exact Bellman algorithm, and---in the smooth positive-density case---the asymptotic value

\[
\boxed{
e_n\sim
\frac1{4n}
\left(\int\sqrt p\right)^2.
}
\]
